%\begingroup
%\let\cleardoublepage\clearpage


% English abstract
\cleardoublepage
\chapter*{Abstract}
\markboth{Abstract}{Abstract}
\addcontentsline{toc}{chapter}{Abstract} % adds an entry to the table of contents
This thesis investigates the origin of the persistent decline of the inner density in a collisionless dark matter halo evolved in isolation for 100 Gyr with the SWIFT code. The central question is whether the apparent cusp-to-core evolution is driven by genuine collisionless dynamics, by numerical artefacts in the gravity solver, or by post-processing choices. To address this, the thesis combines a baseline analysis of the density evolution with a set of controlled tests on the simulation volume, the initial halo structure, the gravitational softening length, and the self-gravity solver parameters. Additional diagnostics based on cumulative mass, shell decomposition, and kinematics are used to distinguish mass redistribution from numerical heating.\\

The results show that the inner density declines steadily throughout the full integration time, but the effect is not removed by changes in box size, softening, or solver settings. A cored initial condition still evolves toward lower central density, although with reduced amplitude, and the hybrid truncated-NFW experiments indicate that suppressing long range forces only weakly alters the trend until the cutoff approaches the softening scale. The kinematic evolution is consistent with numerical heating, the central velocity dispersion rises while the anisotropy becomes more tangential. Taken together, these findings rule out a simple centering or analysis artefact and point instead to residual numerical heating associated with finite-$N$ relaxation and discreteness noise as the most likely cause of the observed core evolution.


% French abstract
\begin{otherlanguage}{french}
\cleardoublepage
\chapter*{Résumé}
\markboth{Résumé}{Résumé}
Cette thèse examine l'origine de la diminution persistante de la densité interne dans un halo de matière noire non-collisionnelle évoluant de manière isolée pendant 100 Gyr avec le code SWIFT. La question centrale est de savoir si l'évolution apparente d'un profil à pointe vers un cœur provient d'une dynamique non-collisionnelle réelle, d'artefacts numériques dans le solveur gravitationnel ou de choix de post-traitement. Pour y répondre, la thèse combine une analyse de référence de l'évolution de la densité avec des tests contrôlés portant sur le volume de simulation, la structure initiale du halo, la longueur d'adoucissement gravitationnel et les paramètres du solveur d'auto-gravité. Des diagnostics complémentaires fondés sur la masse cumulée, la décomposition en couches sphériques et la cinématique permettent de distinguer une redistribution de masse d'un chauffage numérique.\\

Les résultats montrent que la densité interne décroît de façon continue tout au long du temps d'intégration, mais que cet effet n'est éliminé ni par la taille de la boîte, ni par l'adoucissement, ni par les paramètres du solveur. Une condition initiale à cœur évolue elle aussi vers une densité centrale plus faible, quoique avec une amplitude réduite, et les expériences hybrides de type NFW tronqué indiquent que la suppression des forces à longue portée ne modifie que faiblement la tendance jusqu'à ce que le rayon de coupure approche l'échelle d'adoucissement. L'évolution cinématique est compatible avec un chauffage numérique, la dispersion de vitesse centrale augmente tandis que l'anisotropie devient plus tangentielle. L'ensemble de ces résultats exclut un simple artefact de centrage ou d'analyse, et indique plutôt un chauffage numérique résiduel, lié à la relaxation à $N$ fini et au bruit de discrétisation, comme cause la plus probable de l'évolution observée du cœur.
\end{otherlanguage}


%\endgroup			
%\vfill
